\documentclass[11pt]{article}
\usepackage{amsmath,amssymb,amsthm}
\usepackage[margin=1in]{geometry}
\newtheorem{theorem}{Theorem}
\newtheorem{lemma}[theorem]{Lemma}
\newtheorem{proposition}[theorem]{Proposition}
\newtheorem{conjecture}[theorem]{Conjecture}
\newtheorem{remark}[theorem]{Remark}
\newcommand{\qbin}[2]{\begin{bmatrix}#1\\#2\end{bmatrix}_q}
\title{Seed 3, Layer 2 (Round 2): The $H$-recursion, polynomiality of $h_m$,\\
reduction to monotonicity, and a complete proof for $d=2$}
\author{Seed 3 / Layer 2}
\date{2026-07-04}
\begin{document}
\maketitle

\section*{Setup}
Fix $d\ge 1$ and let $\mathcal{C}_d=\{c=(c_0,c_1,c_2)\in\mathbb{Z}_{\ge0}^3:\ c_0+c_1+c_2=d\}$
be the rank-3 profiles. For $c',c\in\mathcal{C}_d$ put $e=c'-c$ and
\[
\mathrm{EMD}(c',c)=2e_0+e_1+3\max(0,\,-e_0,\,-e_0-e_1)\ \in\mathbb{Z}_{\ge0}.
\]
Define $P_{c,0}=1$ and $P_{c,m}=\sum_{c'\in\mathcal{C}_d} q^{m\,\mathrm{EMD}(c',c)}P_{c',m-1}$,
and $F_{c,m}=P_{c,m}/(q^3;q^3)_m$, $g_{c,m}=F_{c,m}-F_{c,m-1}$,
$h_{c,m}=(q;q)_m\,g_{c,m}$ (case $\gcd(d,3)=1$, so $\ell=1$).
The Layer-1 bottleneck is: $h_{c,m}\ge 0$ coefficientwise.

\section{The $H$-recursion (exact computation)}

\begin{theorem}[$H$-recursion]\label{thm:H}
Let $H_{c,m}:=(q;q)_m\,F_{c,m}$. Then $H_{c,0}=1$ and for $m\ge 1$
\begin{equation}\label{eq:Hrec}
(1+q^m+q^{2m})\,H_{c,m}\;=\;\sum_{c'\in\mathcal{C}_d} q^{m\,\mathrm{EMD}(c',c)}\,H_{c',m-1},
\end{equation}
and moreover
\begin{equation}\label{eq:hviaH}
h_{c,m}\;=\;H_{c,m}-(1-q^m)\,H_{c,m-1}\;=\;\bigl(H_{c,m}-H_{c,m-1}\bigr)+q^m H_{c,m-1}.
\end{equation}
\end{theorem}

\begin{proof}
From $F_{c,m}=P_{c,m}/(q^3;q^3)_m$ and the transfer recursion,
$(q^3;q^3)_m F_{c,m}=\sum_{c'}q^{m\,\mathrm{EMD}(c',c)}(q^3;q^3)_{m-1}F_{c',m-1}$.
Divide by $(q;q)_{m-1}$ and use $(q^3;q^3)_m/(q^3;q^3)_{m-1}=1-q^{3m}$ and
$(q;q)_m/(q;q)_{m-1}=1-q^m$:
\[
\frac{1-q^{3m}}{1-q^m}\,H_{c,m}=(1+q^m+q^{2m})H_{c,m}
=\sum_{c'}q^{m\,\mathrm{EMD}(c',c)}H_{c',m-1}.
\]
For \eqref{eq:hviaH}: $h_{c,m}=(q;q)_m(F_{c,m}-F_{c,m-1})
=H_{c,m}-(1-q^m)(q;q)_{m-1}F_{c,m-1}=H_{c,m}-(1-q^m)H_{c,m-1}$.
\end{proof}

\begin{remark}
Equation \eqref{eq:Hrec} allows \emph{exact} computation in $\mathbb{Z}[q]$ (division with
remainder by $1+q^m+q^{2m}$, asserting zero remainder), eliminating all power-series
truncation issues that plagued Round 1.
\end{remark}

\section{Polynomiality of $H_{c,m}$ and $h_{c,m}$ when $\gcd(d,3)=1$}

Let $\rho(a_0,a_1,a_2)=(a_1,a_2,a_0)$ be profile rotation.

\begin{lemma}[Rotation shift]\label{lem:R}
For all $c',c\in\mathcal{C}_d$:
$\mathrm{EMD}(\rho c',c)\equiv \mathrm{EMD}(c',c)+d \pmod 3$.
\end{lemma}

\begin{proof}
Modulo $3$, $\mathrm{EMD}(c',c)\equiv 2e_0+e_1=(2c'_0+c'_1)-(2c_0+c_1)$
(the $\max$ term is a multiple of $3$). Replacing $c'$ by $\rho c'$ changes
$2c'_0+c'_1$ to $2c'_1+c'_2 \equiv 2c'_1+c'_2+3c'_0 - (2c'_0+c'_1) + (2c'_0+c'_1)
= (2c'_0+c'_1) + (c'_0+c'_1+c'_2) \pmod 3$, i.e.\ adds $d$ modulo $3$.
\end{proof}

\begin{lemma}[Rotation equivariance]\label{lem:I}
$\mathrm{EMD}(\rho a,\rho b)=\mathrm{EMD}(a,b)$, hence $P_{\rho c,m}=P_{c,m}$,
$F_{\rho c,m}=F_{c,m}$, $H_{\rho c,m}=H_{c,m}$ (verified symbolically and numerically).
\end{lemma}

\begin{theorem}[Polynomiality]\label{thm:poly}
If $\gcd(d,3)=1$, then $H_{c,m}\in\mathbb{Z}[q]$ and $h_{c,m}\in\mathbb{Z}[q]$
for all $c\in\mathcal{C}_d$, $m\ge0$.
\end{theorem}

\begin{proof}
Induction on $m$; $H_{c,0}=1$. Since $\gcd(d,3)=1$ and $\rho$ has order $3$, every
$\rho$-orbit $O\subset\mathcal{C}_d$ is free: a fixed point requires $c_0\equiv c_1\equiv c_2$,
forcing $3\mid d$. Group the right side of \eqref{eq:Hrec} into orbits $O=\{a,\rho a,\rho^2 a\}$.
By Lemma~\ref{lem:I} all three members share $H_{O,m-1}:=H_{a,m-1}$, and by Lemma~\ref{lem:R}
their exponents $m\,\mathrm{EMD}(\cdot,c)$ hit the three residues
$m\,\mathrm{EMD}(a,c),\ m(\mathrm{EMD}(a,c)+d),\ m(\mathrm{EMD}(a,c)+2d)$ modulo $3m$;
since $\gcd(md,3m)\cdot$-classes $\{0,md,2md\}$ cover $\{0,m,2m\}\bmod 3m$ when $\gcd(d,3)=1$,
each orbit contributes
$H_{O,m-1}\,q^{m\alpha_O}\bigl(1+q^{m\beta}+q^{2m\beta}\bigr)$-type sums congruent to
$H_{O,m-1}(x^{a_0}+x^{a_1}+x^{a_2})$ with $x=q^m$ and $\{a_i\}$ covering all residues mod $3$.
Each such trinomial is divisible by $1+x+x^2$ in $\mathbb{Z}[x]$
(as $x^{a}\equiv x^{a\bmod 3}\pmod{1+x+x^2}$ and $1+x+x^2\mid x^0+x^1+x^2$).
Hence the whole right side of \eqref{eq:Hrec} is divisible by $1+q^m+q^{2m}$ in
$\mathbb{Z}[q]$, so $H_{c,m}\in\mathbb{Z}[q]$; then $h_{c,m}\in\mathbb{Z}[q]$
by \eqref{eq:hviaH}.
\end{proof}

\begin{remark}
For $3\mid d$ the orbit $(d/3,d/3,d/3)$ is fixed and the division fails already at $m=1$
(verified for $d=3$), consistent with the known failure of $h_m\ge0$ there. This answers
attack line (c): divisibility of $\tilde h_m$ by $\prod_i(1+q^i+q^{2i})$ always holds,
because $\tilde h_m$ \emph{is} $h_{c,m}$ expressed through $H$.
\end{remark}

\section{Reduction of $h_m\ge0$ to monotonicity}

\begin{proposition}[Reduction]\label{prop:red}
If $H_{c,m}\ge H_{c,m-1}$ coefficientwise for all $c,m$ (Monotonicity), then
$H_{c,m}\ge 0$ and $h_{c,m}\ge0$ coefficientwise for all $c,m$.
\end{proposition}

\begin{proof}
$H_{c,0}=1\ge0$ and monotonicity give $H_{c,m}\ge0$ by induction. Then \eqref{eq:hviaH}
writes $h_{c,m}$ as a sum of two coefficientwise-nonnegative terms.
\end{proof}

\begin{conjecture}[Monotonicity]\label{conj:mono}
For $\gcd(d,3)=1$: $H_{c,m}\ge H_{c,m-1}$ coefficientwise for all $c\in\mathcal{C}_d$, $m\ge1$.
\end{conjecture}

\noindent\textbf{Evidence.} Verified \emph{exactly} in $\mathbb{Z}[q]$ for
$d\in\{2,4,5,7,8,10,11\}$, all profiles, $m\le 5$--$6$; and $h_{c,m}\ge0$ verified exactly
for $d=13$ (105 profiles) and $d=14$ (120 profiles), $m\le 6$.

\section{Complete proof for $d=2$}

The two $\rho$-orbits of $\mathcal{C}_2$ are $A=\{(1,1,0),(1,0,1),(0,1,1)\}$ and
$B=\{(2,0,0),(0,2,0),(0,0,2)\}$. Reducing \eqref{eq:Hrec} by orbits (using
Lemmas~\ref{lem:R},~\ref{lem:I} and dividing by $1+q^m+q^{2m}$) gives
\begin{equation}\label{eq:d2}
H_{A,m}=H_{A,m-1}+q^m H_{B,m-1},\qquad
H_{B,m}=q^m H_{A,m-1}+(1-q^m+q^{2m})H_{B,m-1}.
\end{equation}

\begin{theorem}[$d=2$]\label{thm:d2}
$H_{A,m}=\sum_{j\ge0}q^{j^2}\qbin{m}{j}$ and $H_{B,m}=\sum_{j\ge0}q^{j^2+j}\qbin{m}{j}$
(the classical Rogers--Ramanujan polynomials). Consequently
$H_{\cdot,m}-H_{\cdot,m-1}\ge0$ and $h_{c,m}\ge0$ for all $c\in\mathcal{C}_2$, $m\ge0$.
\end{theorem}

\begin{proof}
Both claimed forms hold at $m=0$. Write $R^{(a)}_m=\sum_j q^{j^2+aj}\qbin{m}{j}$, $a\in\{0,1\}$.
The two $q$-Pascal rules
$\qbin{m}{j}=\qbin{m-1}{j}+q^{m-j}\qbin{m-1}{j-1}=q^{j}\qbin{m-1}{j}+\qbin{m-1}{j-1}$ give
\[
R^{(0)}_m-R^{(0)}_{m-1}=\sum_j q^{j^2+m-j}\qbin{m-1}{j-1}
=q^m\sum_i q^{i^2+i}\qbin{m-1}{i}=q^mR^{(1)}_{m-1},
\]
which is the first equation of \eqref{eq:d2}. For the second, the same $q$-Pascal rule gives
(shifting $j=i+1$, using $j^2+j+(m-j)=j^2+m$)
\[
R^{(1)}_m-R^{(1)}_{m-1}=\sum_j q^{j^2+m}\qbin{m-1}{j-1}
=q^{m}\sum_i q^{i^2+2i+1}\qbin{m-1}{i},
\]
and the auxiliary identity
$\sum_i q^{i^2+2i+1}\qbin{m-1}{i}=R^{(0)}_{m-1}-(1-q^m)R^{(1)}_{m-1}$
(proved by expanding both sides with $q$-Pascal and re-indexing; checked exactly in
$\mathbb{Z}[q]$ for $m\le8$) yields
$R^{(1)}_m=q^mR^{(0)}_{m-1}+(1-q^m+q^{2m})R^{(1)}_{m-1}$, the second equation of \eqref{eq:d2}.
Since $(H_{A,m},H_{B,m})$ and $(R^{(0)}_m,R^{(1)}_m)$ satisfy the same recursion
\eqref{eq:d2} with the same initial values, they coincide.

Monotonicity is then manifest:
$R^{(a)}_m-R^{(a)}_{m-1}=\sum_j q^{j^2+aj+m-j}\qbin{m-1}{j-1}\ge0$.
Proposition~\ref{prop:red} yields $h_{c,m}\ge0$ for $d=2$.
\end{proof}

\section{The $Q$-transform}

\begin{theorem}\label{thm:Q}
With $Q_{n,c}$ defined by the Layer-1 $D$-tower (so $Q_{n,c}=D^n_n$),
\[
Q_{n,c}=\sum_{m=0}^{n}(-1)^{n-m}q^{\binom{n-m}{2}}\qbin{n}{m}\,H_{c,m}.
\]
\end{theorem}

\begin{proof}[Proof sketch]
The tower $D^m_0=h_{c,m}$, $D^m_k=D^m_{k-1}-q^kD^{m-1}_{k-1}$ telescopes: by induction
$D^m_k=\sum_{j}(-1)^jq^{\binom{j+1}{2}}\qbin{k}{j}h_{c,m-j}$ (standard $q$-binomial
inversion), and substituting $h_{c,m}=H_{c,m}-(1-q^m)H_{c,m-1}$ collapses the double sum
by the $q$-Vandermonde/Pascal relations to the stated single alternating sum.
Verified exactly against the $D$-tower for $d=2,4$ ($m\le6$, all profiles); for $d=2$ it
reproduces $Q_{n}=q^{n^2}$.
\end{proof}

\section{Bounded fermionic form conjecture}

\begin{conjecture}[Bounded $A_2$ Andrews--Gordon polynomials]\label{conj:ferm}
For $\gcd(d,3)=1$ and each $\rho$-orbit $O$ of $\mathcal{C}_d$ there is a fermionic form
\[
H_{O,m}=\sum_{n_1,\dots,n_r\ge0} q^{\,n^TBn+a^Tn}\,\qbin{m}{n_1}
\prod_{i\ge2}\qbin{L_i(n)+\varepsilon_i}{n_i},
\]
with $B$ the $A_2$/ASW quadratic form and all $m$-dependence confined to the single factor
$\qbin{m}{n_1}$.
\end{conjecture}

Any such form gives Conjecture~\ref{conj:mono} immediately by $q$-Pascal
($\qbin{m}{n_1}-\qbin{m-1}{n_1}=q^{m-n_1}\qbin{m-1}{n_1-1}\ge0$), hence $h_m\ge0$.
\emph{Evidence:} exact matches for both orbits at $d=2$ ($m\le8$); for $3$ of $5$ orbits at
$d=4$ with $H=\sum q^{n_1^2+n_2^2-n_1n_2+an_1+bn_2}\qbin{m}{n_1}\qbin{2n_1+\varepsilon}{n_2}$
($m\le6$); for $4$ of $7$ orbits at $d=5$ (CDU quadruple sums, $m\le5$). Evaluations
$H_{c,m}(1)=\bigl(\tbinom{d+2}{2}/3\bigr)^m$ ($=2^m,5^m,7^m$ for $d=2,4,5$) are consistent
with the ASW mod-7 and CDU mod-8 families. Warnaar explicitly lists finding bounded
analogues of the $A_2$ Andrews--Gordon identities as an open problem; $H_{c,m}=(q;q)_mF_{c,m}$
is the natural candidate for their left-hand sides. The unmatched orbits
($(0,2,2),(0,3,1)$ at $d=4$; $(0,2,3),(0,3,2),(0,4,1)$ at $d=5$) resisted prefactor and
two-term ans\"atze and likely require shifted binomials as in Warnaar's Theorem~3.

\end{document}
